% ==================================================
% MODELO VAZIO — ALTITUDE CENTRO UNIVERSITÁRIO
% Preencha somente os dados e o conteúdo do curso.
% Logo, cabeçalho, rodapé e identidade visual são aplicados pelo portal.
% ==================================================

\documentclass{article}

% ==================================================
% DADOS DO CURSO
% ==================================================

\codigocurso{CODIGO-DO-CURSO}
\titulocurso{Título do curso}
\areacurso{Área do curso}
\nivelcurso{BASICO}
\notaminima{70}
\descricaoCurso{Escreva uma descrição breve do curso.}

\begin{document}

% ==================================================
% MÓDULO 1
% Repita o ambiente altitudemodulo para criar outros módulos.
% ==================================================

\begin{altitudemodulo}{Título do módulo}{1}
\descricaoModulo{Descrição breve do módulo.}

% Vídeo opcional do módulo. Apague a linha quando não houver vídeo.
% \videoModulo{https://www.youtube.com/watch?v=SEU_VIDEO}

\begin{conteudo}

\section{Título da seção}

Digite o desenvolvimento do conteúdo aqui.

\subsection{Título da subseção}

Continue o conteúdo aqui.

\begin{itemize}
    \item Primeiro item;
    \item Segundo item.
\end{itemize}

\end{conteudo}

% ==================================================
% QUESTÃO — repita o bloco para adicionar outras questões
% ==================================================

\begin{altitudequestao}{1}
\enunciado{Digite o enunciado da questão.}
\alternativa{A}{Alternativa A.}
\alternativa{B}{Alternativa B.}
\alternativa{C}{Alternativa C.}
\alternativa{D}{Alternativa D.}
\alternativa{E}{Alternativa E.}
\gabarito{A}
\resolucao{Explique de forma objetiva por que a alternativa está correta.}
\end{altitudequestao}

\end{altitudemodulo}

\end{document}
